\subsection{The rigid description}
\label{subsec:rigid}

In the rigid model the transition of $X_9$ has two descriptions. At the four-node
roots it is the hyperk\"ahler quotient of four hypermultiplets by a rank-three torus. At the
$E_2$ point it is built from the two nodal hypermultiplets and the $E_2$ theory, whose reduced
Higgs branch is the cone $\CC[J_+,J_0,J_-]/(J_+J_--J_0^2)=\CC^2/\ZZ_2$ generated by the moment maps
$J_\pm,J_0$ of its $SU(2)$ flavour symmetry; the two bulk vector multiplets gauge a rank-two
torus acting on the nodes and on this cone~\cite[eqs.~(2.3), (3.10)]{W26quantum}. Both give
$\CC[X,Y,s]/(XY-s^4)$, the sign of $s^4$ being a convention (compare~\eqref{eq:app-ring}).
The two presentations are related by reduction in stages: quotienting the four-hypermultiplet
model first by the third row of~\eqref{eq:app-Q} produces the $\CC^2/\ZZ_2$ cone of the $E_2$
theory, on which the first two rows act as the rank-two gauging.
Theorems~\ref{thm:four-node} and~\ref{thm:wall} show that this rigid germ survives at finite
Planck mass under Hypothesis~\ref{hyp:regular}, at the four-node roots and, for the completion of
the regular Higgs branch, at the del Pezzo facet.

\subsection{The other roots of \texorpdfstring{$X_9$}{X9}}
\label{subsec:other-roots}

Theorem~\ref{thm:wall} applies over the whole region~\eqref{eq:smooth-kahler}, but the root over
which the regular Higgs vacua end changes across the region. Table~\ref{tab:roots} lists the
roots over $y_3,y_4>0$. The regular Higgs vacua over $y_2>0$ end at the four-node roots. Those
over $\max(-y_3,-2y_4)<y_2<0$ end at the locus where the image of $E$ in the flop of $\wideX_9$
along the curve $h-e_1-e_2$, a copy of $\mathbb P^1\times\mathbb P^1$, collapses to a point, giving the rank-one $E_1$ theory together with the two nodes. Across the wall $y_2=-2y_4$ the toric divisor
$D_2$ of Section~\ref{subsec:x9} collapses together with $C_1$ and $C_2$, and beyond it two surfaces are collapsed at the root. The
cubic prepotential is continuous across each of these walls, as an identity among intersection
numbers. The lists of curves of zero area rest on the Gopakumar--Vafa invariants of all classes
of total degree at most eight; for curves $c=dh-d_1e_1-d_2e_2$ inside $E$, whose area at the
$E_1$ root is $\tfrac{y_2}2(d_1+d_2-d)$, they hold in all degrees.

\begin{table}[t]
\centering\small
\begin{tabular}{lll}
\toprule
Higgs-side region & Coulomb-side root & light states at the root\\
\midrule
$y_2>0$ & $J=y\in F$ & four free hypermultiplets\\
$y_2=0$ & $J=y\in F_0$ & $E_2$ theory and two nodes\\
$\max(-y_3,-2y_4)<y_2<0$ & $J=y+\tfrac{y_2}2E$ (after a flop) & $E_1$ theory and two nodes\\
$y_2=-2y_4>-y_3$ & wall & $D_2\cup C_1\cup C_2$ contracts to a point\\
$-y_3<y_2<-2y_4$ & $J=y+\tfrac{y_2}2E+\tfrac{y_2+2y_4}2D_2$ & two collapsed surfaces\\
\bottomrule
\end{tabular}
\caption{Roots of the Higgs branch over the region $y_3,y_4>0$ of the K\"ahler cone of
$X_{9,t}$, with $y=y_2H_2+y_3H_3+y_4H_4$. At $y_2=-y_3$ a plane in $X_{9,t}$ collapses and the
vector-multiplet metric on the Higgs side degenerates.}
\label{tab:roots}
\end{table}

These rows are not used in the proofs. At all these roots Theorem~\ref{thm:wall} describes the
completion of the regular Higgs branch; whether further Higgs-type components meet the Coulomb
branch there has not been examined. Theorem~\ref{thm:wall} does not apply beyond $y_2=-y_3$. There
a divisor isomorphic to $\mathbb P^2$ in $X_{9,t}$ collapses, giving the $E_0$ theory on the
smooth side~\cite{W26quantum}, the Hessian
\begin{equation}
 \det\operatorname{Hess}\bigl(\kappa|_{X_{9,t}}\bigr)\;\propto\;
 (y_2+y_3)(y_2+y_3+3y_4)(y_2+2y_3+4y_4)
 \label{eq:hess}
\end{equation}
of the cubic form of $X_{9,t}$ vanishes, and the regular Higgs branch itself degenerates.

The facet $y_3=0$ of $\overline F$, which we call the $SU(2)$ facet because the curve class
$\beta_3$ becomes massless there and gives an $SU(2)$ vector multiplet, provides a consistency
test of the definition of the Higgs branch as a metric completion, at a root with a Lagrangian description that lies outside the scope of
Theorem~\ref{thm:normal-form}. If no class of degree $(0,j,0)$, $j\ge1$, is effective on
$X_{9,t}$ (for $j=1,2$ the Gopakumar--Vafa invariants in these degrees sum to zero), the relative
interior of this facet is interior to $\mathcal K(X_{9,t})$, and the completion there again has
hypermultiplet germ $N\times\CC^2/\ZZ_4$. Among classes of total degree at most eight, the light
states on the Coulomb side of this facet are an $SU(2)$ vector multiplet from the class
$\beta_3$, two doublets $(e_i,e_i+\beta_3)$ and the singlets $C_1,C_2$; the $SU(2)$ coupling is
finite, since it is set by the area $y_2>0$ of $h-e_1-e_2$. The rigid model of this
$SU(2)\times U(1)^3$ theory is again exactly $\CC^2/\ZZ_4$. The level set of the complex moment
map has a single component of stable points, on which the two doublets are aligned, and the
invariant ring of the whole level set is $\CC[A,B,s]/(AB+s^4)$. This agrees with the germ of the completion, but it
concerns the rigid model only: the generic stabiliser there is one-dimensional, so
Theorem~\ref{thm:normal-form} does not apply at finite Planck mass.

\subsection{Instanton geometry without gravity}
\label{subsec:instantons}

In type~IIA string theory on a threefold near a conifold degeneration, the hypermultiplet
moduli space receives corrections from Euclidean D2-branes on the vanishing three-sphere, and
the corrected geometry near $k$ homologous vanishing three-cycles has a $\CC^2/\ZZ_k$
singularity~\cite{OoguriVafa:1996,SeibergShenker:1996}; the multiplicity is the
Gopakumar--Vafa invariant of the vanishing class~\cite{SaueressigVandoren:2007}, and the smooth
model geometries of this kind were constructed rigorously for $k=1$ in~\cite{FredricksonZimet:2025}, where the extension to multiplicity $k$ is
indicated~\cite[Rems.~3.14, 4.44]{FredricksonZimet:2025}. For $X_9$ the
relation~\eqref{eq:face-relation} has unit coefficients, so the four vanishing three-spheres of
the smoothing of $X_F$ lie in the classes $\pm[S]$ of one primitive class, and this picture gives
an $A_3$ point, in agreement with Theorem~\ref{thm:four-node}. These analyses rest on rigid,
gravity-decoupled arguments. Greene, Morrison and Vafa derived $\CC^2/\ZZ_k$ for $k$
hypermultiplets with one relation of this kind~\cite[Sect.~3]{Greene:1996Confinement} and
expected that with gravity the qualitative features of the four-dimensional hypermultiplet space
persist while the metric is corrected~\cite[Sect.~6]{Greene:1996Confinement}.
Theorem~\ref{thm:general} and Proposition~\ref{prop:tangent-cone} make the five-dimensional
counterpart precise, under the regularity hypothesis, for every face satisfying (S1)--(S3),
$X_9$ included: the germ is unchanged and the metric changes only at relative order $r^2$.
Corollary~\ref{cor:gmv} gives the four-dimensional statement about the type~IIA hypermultiplet moduli
space, for the smoothings of the contractions $X_F$ of faces satisfying (S1)--(S3) and under the identification (W-IIA),
and Proposition~\ref{prop:x9-leading} identifies the leading correction to the Wilsonian quotient
metric, which is the moduli-space metric if (W-IIA) holds isometrically.
Theorem~\ref{thm:general} does not cover faces whose curves satisfy several relations; there
proper subsets of the light states can condense, the flat quotient has non-isolated
singularities, and Theorem~\ref{thm:normal-form} would need a stratified extension.

\subsection{A census of one-relation faces}
\label{subsec:census}

Theorem~\ref{thm:general} applies to every face with one relation, and $X_9$ is one example. Its proof uses
the light spectrum, the openness of $F$ in the annihilator of the classes of the light states, and the
positivity of the kinetic matrix. The census below tests these three inputs and, where toric data suffice,
the geometric part of (S1), namely that the light states come from disjoint $(-1,-1)$-curves; it does not
test the saturation in (S2), which only normalises the
torus. To see how often these hold, and whether coefficients $|a_i|\geq2$ occur,
we examined all favourable reflexive polytopes (those for which every divisor class of the hypersurface
is the restriction of an ambient toric divisor class) in four dimensions with $2\leq h^{1,1}\leq5$ in the
classification of Kreuzer and Skarke~\cite{Kreuzer:2000}, and a random sample of $2000$ of the $17101$
polytopes with $h^{1,1}=6$, of which the favourable ones were examined. For each polytope we took all fine regular star triangulations when there are
at most forty, and otherwise up to forty distinct ones obtained from random heights. For each triangulation we took the
closed cone of K\"ahler classes of the hypersurface pulled back from the toric ambient space, which we call
its toric K\"ahler cone, and every proper face $F$ of this cone. At an interior point of $F$ we listed the
curve classes of zero area with nonzero genus-zero Gopakumar--Vafa invariant $n_\beta$, up to degree six
with respect to a grading vector positive on the Mori cone of the toric ambient space, computed from the mirror
periods~\cite{Demirtas:2024Mirror} with the software of~\cite{Demirtas:2022cytools}. We call $F$ a \emph{one-relation face} when
\begin{enumerate}
\item[(i)] all these invariants are positive;
\item[(ii)] the volume and the volumes of all prime toric divisors are positive, computed exactly from the
 intersection numbers; together with (i) this excludes massless vector multiplets (invariant $-2$) and
 prime toric divisors that are ruled surfaces collapsing onto curves of genus $g\geq2$, whose fibre class
 has the positive invariant $2g-2$;
\item[(iii)] the light spectrum, with $n_\beta$ hypermultiplets of charge $\beta$ and $k=\sum_\beta n_\beta$,
 has exactly one relation among its charges, with all coefficients nonzero (so $k\geq2$);
\item[(iv)] $F$ is open in the annihilator of these classes, as in (S1);
\item[(v)] the gauge kinetic matrix is positive definite at that point.
\end{enumerate}
A face found from several triangulations of one polytope, with the same classes up to sign, was counted
once. Table~\ref{tab:census} gives the result.

\begin{table}[t]
\centering\small
\begin{tabular}{rrrrrrr}
\toprule
$h^{1,1}$ & scanned & $k=2$ & $k=3$ & $k=4$ & $k=5$ & $|a_i|\geq2$\\
\midrule
2 & 36 & 1 (1) & -- & -- & -- & 0\\
3 & 243 & 24 (21) & 1 (1) & -- & -- & 0\\
4 & 1185 & 335 (243) & 11 (8) & -- & -- & 0\\
5 & 4897 & 2895 (1727) & 83 (56) & 113 (63) & 1 (1) & 0\\
6 & 1948 & 1978 (946) & 54 (33) & 108 (52) & 1 (1) & 0\\
\bottomrule
\end{tabular}
\caption{One-relation faces, faces satisfying (i)--(v), by $h^{1,1}$ and the number $k=\sum_\beta n_\beta$ of
light hypermultiplets, in Calabi--Yau hypersurfaces from favourable Kreuzer--Skarke polytopes (all
polytopes with $h^{1,1}\leq5$; the favourable ones among a random sample of $2000$ of the polytopes with
$h^{1,1}=6$). ``Scanned'' is
the number of favourable polytopes examined. Entries give the number of faces and, in parentheses, of
polytopes; a polytope can appear in several columns. All coefficients satisfy $|a_i|=1$, so $m=k$.}
\label{tab:census}
\end{table}

Most one-relation faces are flop walls, $k=2$, where two hypermultiplets of equal charge become massless, the
situation of the gauged supergravities of~\cite{Mohaupt:2004de}; faces with $k=3$, $4$ or $5$ occur in
210
 polytopes. In no case does a coefficient $|a_i|\geq2$ occur. For a flop wall, on which all classes
of zero area lie on one ray, this would require exactly one isolated curve in each of two classes
$p\beta_0$ and $p'\beta_0$, $p<p'$, and none in the other multiples, for example $\beta$ and $2\beta$; the
relation is then $a=(p',-p)/\gcd(p,p')$. Here the length $d$ of a flopping curve is that of Katz and Morrison~\cite{KatzMorrison:1992}, the
coefficient of the curve in the fundamental cycle of a general hyperplane section through the singular
point. A generic deformation splits a flopping curve of length $d$ into
$n_{j\beta}$ isolated curves in the classes $j\beta$, $j\leq d$~\cite{Katz:2006contractible}, and for the
flopping curves of length two in~\cite{BrownWemyss:2018} one has $(n_\beta,n_{2\beta})=(5,1)$, so this pattern
does not arise from them. Whether faces with $m\neq k$ occur at all remains open;
Theorem~\ref{thm:general} covers them if they do.

Genus-zero invariants alone do not distinguish $k$ disjoint $(-1,-1)$-curves, as for $X_9$, from other
configurations with the same light spectrum, for instance two curves meeting in a point, contracted to a
compound Du Val point of type $cA_2$ (for $k=3$), or a single curve with $n_\beta=2$ in place of two curves
in one class. The geometric reading of Theorem~\ref{thm:general} as a transition through a threefold with
$k$ nodes needs the first alternative. We tested it with toric curves. For a two-dimensional cone
$\{i,j\}$ of the fan, $C_{ij}=D_i\cdot D_j\cdot X$ is, for generic $X$, empty or a disjoint union of $\ell$
torus-translates of a smooth curve of genus $g$, with $\ell>1$ only if $g=0$; its class is given by the
intersection numbers $\kappa(D,D_i,D_j)$, and by adjunction $\kappa_{iij}+\kappa_{ijj}=\ell(2g-2)$, so normal
degrees $(\kappa_{iij},\kappa_{ijj})=(-1,-1)$ force a single $(-1,-1)$-curve and $(-2,-2)$ two disjoint ones
in half the class. Two such curves with one common divisor meet in $\kappa_{ijl}$ points, and curves without
a common divisor are disjoint, since $X$ avoids the torus-fixed points. If every ray of the fan is a ray of
the fan of the ambient toric morphism defined by $F$ (equivalently, $J^3\cdot D_i>0$ on the ambient space for
every prime toric divisor $D_i$), every curve of zero area is a component of some $C_{ij}$; this condition
depends on $F$ alone. We say that toric data \emph{decide} a face satisfying it if, in some triangulation whose toric K\"ahler cone
contains $F$, the zero-area curves $C_{ij}$ are exactly the light states, with normal degrees $(-1,-1)$ (for
a light class with $n_\beta=2$, either two such curves or one $C_{ij}$ in twice the class with degrees
$(-2,-2)$), and are pairwise disjoint. All such triangulations give small resolutions of the same
contraction $X_F$, so a decision in one of them shows that $X_F$ has $k$ nodes. Toric data decide the four-node face of
$X_9$ in this way. Every triangulation containing a face was examined (all of them when there
are at most 40, otherwise a sample, as in the census). Of the 372 faces with $k\geq3$, toric data decide 335: all
with $k=5$, 212 of the 221 with $k=4$ and 121 of the 149 with
$k=3$. In the others a toric divisor is contracted, so curves of zero area need not be toric: in the
remaining 9 faces with $k=4$ the light states are toric, pairwise disjoint $(-1,-1)$-curves, but further
curves of zero area are not excluded by toric data, and in the remaining 28 faces with $k=3$ at least one
light class (exactly one in 24 of them) has no toric representative in any triangulation containing the
face. Of the 5233 flop walls, toric data decide 3757; in the other 1476 a toric
divisor is contracted and the light class has no toric representative in the triangulations examined, so
these are flops of non-toric curves. In all 5605 faces and all triangulations examined, no two zero-area
toric curves meet, and every zero-area toric curve has normal degrees $(-1,-1)$ or $(-2,-2)$. Faces of the
K\"ahler cone of the hypersurface that are not faces of a toric K\"ahler cone are not seen.

\subsection{Other effects at finite Planck mass}
\label{subsec:other-effects}

Membrane instantons on three-cycles that do not vanish, and fivebrane instantons, correct the
metric of the hypermultiplet manifold smoothly and preserve the gauged isometries, so they are
covered by Hypothesis~\ref{hyp:regular}: the moment map of the corrected metric has no constants
of integration and still vanishes at the root, since the Coulomb branch consists of
supersymmetric vacua. Membrane instantons on the vanishing three-spheres are expected, as
in~\cite{OoguriVafa:1996,SeibergShenker:1996}, not to be an additional correction but the
description of the same $A_3$ point in terms of the neutral fields alone. The gauged torus is
exact and its weights are quantised, and the charge lattice is saturated, so no discrete gauge symmetry survives at a generic Higgs vacuum.

The argument of Section~\ref{sec:higgs} is five-dimensional and does not use a compactification on
a circle. At finite radius the four-node face corresponds to the stratum of the four-dimensional
Coulomb branch found in~\cite[Sects.~5.1--5.2]{W26quantum}, on which the same four hypermultiplets
become massless; Section~\ref{sec:4d} treats type~IIA string theory on $\wideX_9$ there directly, where
these states carry no Kaluza--Klein charge.

\subsection{Open problems}
\label{subsec:open}

We list the questions left open.
\begin{enumerate}
\item Is the metric completion at the $E_2$ point the Higgs branch of the interacting $E_2$ theory
 coupled to supergravity? No Lagrangian is available there; the rigid computations at the del Pezzo
 facet (Section~\ref{subsec:rigid}) and, at another root, at the $SU(2)$ facet
 (Section~\ref{subsec:other-roots}) are consistent with defining the Higgs branch as this completion.
\item What is the Higgs branch at finite Planck mass at the $SU(2)$ facet? There the completion of the
 regular Higgs branch is again $N\times\CC^2/\ZZ_4$ (conditionally, Section~\ref{subsec:other-roots}),
 but Theorem~\ref{thm:normal-form} does not apply and only the rigid model has been computed.
\item Is the metric in the uniformising chart of Definition~\ref{def:leading} twice differentiable, or
 smooth, at the vertex, so that the transverse slice is a Riemannian orbifold there and $R_0$ its
 curvature?
\item Is there a complex-analytic version of the local model, on the twistor space or the Swann bundle?
 For negative scalar curvature this needs a slice theorem for pseudo-hyperk\"ahler cones, which the known
 local models~\cite{Mayrand:2022} do not provide.
\item What are the local structure of the quotient and the tangent cones of its length metric at points
 whose stabiliser is a proper subgroup of $G$, in particular along the locally K\"ahler and totally real
 pieces of~\cite[Sect.~5]{DancerSwann:1997}?
\end{enumerate}
